\chapter{Integration und Bedienung}
\label{ch:integration}

Die \cref{ch:objekterkennung,ch:robotersteuerung,ch:sichtpruefung} haben die Module einzeln
beschrieben. Dieses Kapitel zeigt, wie sie zu einem lauffähigen Gesamtsystem zusammengesetzt
werden: wie ein Arbeitszyklus abläuft, wie das System konfiguriert und bedient wird, wie es
auf Fehler reagiert und wie die Ergebnisse protokolliert werden.

% ------------------------------------------------------------------------------
\section{Ablauf eines Zyklus}
\label{sec:int-zyklus}

Ein Arbeitszyklus folgt dem Zustandsautomaten aus \cref{sec:architektur}.
\Cref{fig:sequenz} zeigt die Aufrufe, die der Orchestrator dabei an die Module richtet.
Zuerst liefert die Objekterkennung die Pose des Bauteils. Die Robotersteuerung greift das
Bauteil, legt es an der Prüfposition ab und fährt anschließend in die Ruheposition, damit der
Roboterarm das Sichtfeld der Seitenkamera freigibt. Danach prüft die Sichtprüfung das Bauteil.
Abhängig vom Ergebnis nimmt der Roboter das Bauteil wieder auf, legt es auf der Ablage für
Gut- oder Schlechtteile ab und kehrt in die Ruheposition zurück. Damit ist der Zyklus
beendet, und der nächste kann beginnen.

\begin{figure}[htbp]
  \centering
  \begin{tikzpicture}[
      head/.style={draw, rounded corners, fill=black!5, minimum width=24mm, minimum height=8mm, font=\small},
      msg/.style={-{Stealth[length=2mm]}, thick},
      ret/.style={-{Stealth[length=2mm]}, dashed},
      lbl/.style={font=\scriptsize, above, inner sep=1pt}]
    \def\yend{-11.2}
    \foreach \x/\n/\t in {0/o/Orchestrator, 3.8/m1/M1 Erkennung, 7.6/m2/M2 Roboter, 11.4/m3/M3 Prüfung} {
      \node[head] (\n) at (\x,0) {\t};
      \draw[black!40] (\x,-0.4) -- (\x,\yend);
    }
    \draw[msg] (0,-1.0) -- node[lbl] {\texttt{detect()}} (3.8,-1.0);
    \draw[ret] (3.8,-1.6) -- node[lbl] {\texttt{ObjectPose}} (0,-1.6);
    \draw[msg] (0,-2.4) -- node[lbl, pos=0.75] {\texttt{pick(pose)}} (7.6,-2.4);
    \draw[msg] (0,-3.2) -- node[lbl, pos=0.75] {\texttt{place(inspection)}} (7.6,-3.2);
    \draw[msg] (0,-4.0) -- node[lbl, pos=0.75] {\texttt{home()}} (7.6,-4.0);
    \draw[msg] (0,-4.8) -- node[lbl, pos=0.85] {\texttt{inspect()}} (11.4,-4.8);
    \draw[ret] (11.4,-5.4) -- node[lbl, pos=0.15] {\texttt{InspectionResult}} (0,-5.4);
    \draw[msg] (0,-6.2) -- node[lbl, pos=0.75] {\texttt{pick(inspection)}} (7.6,-6.2);
    \draw[black!50, rounded corners] (-0.6,-6.7) rectangle (12.4,-9.6);
    \node[font=\scriptsize, anchor=north west, fill=black!5, draw=black!50] at (-0.6,-6.7) {alt};
    \node[font=\scriptsize, anchor=west] at (0.4,-7.05) {[Gutteil]};
    \draw[msg] (0,-7.5) -- node[lbl, pos=0.75] {\texttt{place(bin\_good)}} (7.6,-7.5);
    \draw[black!50, dash pattern=on 2pt off 2pt] (-0.6,-8.1) -- (12.4,-8.1);
    \node[font=\scriptsize, anchor=west] at (0.4,-8.45) {[Schlechtteil]};
    \draw[msg] (0,-8.9) -- node[lbl, pos=0.75] {\texttt{place(bin\_bad)}} (7.6,-8.9);
    \draw[msg] (0,-10.4) -- node[lbl, pos=0.75] {\texttt{home()}} (7.6,-10.4);
  \end{tikzpicture}
  \caption{Aufrufe des Orchestrators an die Module während eines Arbeitszyklus}
  \label{fig:sequenz}
\end{figure}

Findet die Objekterkennung kein Bauteil, endet der Zyklus ohne Fehler, bevor sich der Roboter
bewegt. Im Dauerbetrieb wiederholt der Orchestrator den Zyklus, bis er angehalten wird oder
ein Fehler auftritt.

Da die Module nur über ihre Schnittstellen angesprochen werden, läuft dieser Ablauf
unabhängig davon, welche Implementierung eines Moduls gewählt ist. Derselbe Zyklus lässt sich
daher vollständig gemockt, mit gespeicherten Testfotos und simuliertem Roboter oder an der
realen Anlage ausführen (\cref{sec:teststrategie}).

% ------------------------------------------------------------------------------
\section{Konfiguration}
\label{sec:int-konfiguration}

Alle Werte, die vom konkreten Aufbau abhängen, stehen in einer zentralen
Konfigurationsdatei im YAML-Format. Dazu gehören:
\begin{itemize}
  \item Kameraquellen und Auflösungen,
  \item Größe des Arbeitsbereichs, Wörterbuch und Positionen der ArUco-Marker,
  \item Schwellwerte der Bildverarbeitung für Segmentierung und Prüfmerkmale,
  \item Transformation vom Arbeitsbereich in das Roboter-\ac{ks},
  \item Netzwerkadresse, Werkzeug, Geschwindigkeiten, Arbeitsraumgrenzen und feste Posen des
        Roboters,
  \item Sollwerte und Format der Prüfmerkmale,
  \item die standardmäßig verwendete Implementierung jedes Moduls.
\end{itemize}

Die Module werden über eine gemeinsame Fabrikfunktion erzeugt, die zu einem Modul und einem
Variantennamen die passende Implementierung liefert (\cref{tab:varianten}). Diese Funktion
nutzen sowohl der Start über die Kommandozeile als auch die Bedienoberfläche, sodass beide
Wege dasselbe System aufbauen.

\begin{table}[htbp]
  \centering
  \caption{Wählbare Varianten der Module}
  \label{tab:varianten}
  \begin{tabularx}{\textwidth}{@{}l l X@{}}
    \toprule
    Modul & Variante & Bedeutung \\
    \midrule
    Erkennung & \texttt{mock}   & zufällige Pose, keine Kamera \\
              & \texttt{replay} & Objekterkennung auf gespeicherten Testfotos \\
              & \texttt{camera} & Objekterkennung auf dem Live-Bild der Top-down-Kamera \\
              & \texttt{ai}     & lernbasierter Ansatz aus \cref{ch:kigreifen} \\
    \midrule
    Roboter   & \texttt{mock}   & Bewegungen werden nur protokolliert \\
              & \texttt{sim}    & Robotersteuerung gegen die Nachbildung der Steuerung \\
              & \texttt{real}   & realer Roboter \\
    \midrule
    Prüfung   & \texttt{mock}   & zufälliges Prüfergebnis, keine Kamera \\
              & \texttt{replay} & Sichtprüfung auf gespeicherten Testfotos \\
              & \texttt{camera} & Sichtprüfung auf dem Live-Bild der Seitenkamera \\
    \bottomrule
  \end{tabularx}
\end{table}

% ------------------------------------------------------------------------------
\section{Bedienoberfläche}
\label{sec:int-gui}

Für Inbetriebnahme und Betrieb gibt es eine grafische Bedienoberfläche (\ac{gui}), die mit
dem Framework Qt über dessen Python-Anbindung PySide6 umgesetzt ist. Sie ist in fünf Reiter
gegliedert (\cref{tab:gui-reiter}). Unter den Reitern zeigt ein Protokollfenster laufend die
Meldungen aller Module an.

\begin{table}[htbp]
  \centering
  \caption{Reiter der Bedienoberfläche}
  \label{tab:gui-reiter}
  \begin{tabularx}{\textwidth}{@{}l X@{}}
    \toprule
    Reiter & Funktion \\
    \midrule
    Ablauf & Auswahl der Varianten, einzelner Zyklus oder Dauerbetrieb, Anzeige des aktuellen
             Zustands, der Bilder von Erkennung und Prüfung, einer Draufsicht auf den
             Roboterarbeitsraum sowie einer Tabelle und Zählung der Ergebnisse \\
    M1 Erkennung & Bild aus Kamera, Testfoto oder Datei holen und analysieren, optional
             fortlaufend; Anzeige des Kamerabilds mit erkannten Markern und der entzerrten
             Draufsicht mit Pose \\
    M2 Roboter & Verbinden mit Simulator oder realem Roboter, Status, Greifer öffnen und
             schließen, feste Posen anfahren, aktuelle Roboterpose als feste Pose übernehmen
             (Teachen), Stopp und Freigabe \\
    M3 Prüfung & wie M1 für die Seitenkamera; Anzeige der Prüfbereiche und der erkannten
             Merkmale zum Einstellen der Parameter \\
    Einstellungen & alle Werte der Konfigurationsdatei als bearbeitbarer Baum; Änderungen
             wirken sofort, Speichern schreibt die Datei unter Erhalt der Kommentare \\
    \bottomrule
  \end{tabularx}
\end{table}

\begin{figure}[htbp]
  \centering
  \missingfigure[figwidth=0.9\textwidth]{Screenshot Reiter \enquote{Ablauf} während eines Zyklus}
  \caption{Bedienoberfläche, Reiter \enquote{Ablauf}}
  \label{fig:gui-ablauf}
\end{figure}

Alle Reiter greifen auf eine gemeinsame Sitzung zu, die Konfiguration, Kameras und
Roboterverbindung hält. Eine Kamera oder eine Roboterverbindung wird dadurch nur einmal
geöffnet, auch wenn sie in mehreren Reitern genutzt wird, und Einzeltests in den Modulreitern
und der automatische Ablauf verwenden dieselbe Hardware.

Damit die Oberfläche während langer Roboterbewegungen oder Bildverarbeitungsschritte
bedienbar bleibt, laufen diese Aufgaben in eigenen Hintergrund-Threads. Alle Befehle an den
Roboter werden in einem einzigen Thread strikt nacheinander ausgeführt, sodass sich zwei
Bewegungsbefehle nie überschneiden können. Der Zustandsautomat meldet jeden Zustandswechsel an
die Oberfläche, die den aktuellen Zustand farbig in der Ablaufkette hervorhebt.

% ------------------------------------------------------------------------------
\section{Fehlerbehandlung und Sicherheit im Ablauf}
\label{sec:int-fehler}

Jedes Modul meldet Fehler über eine eigene Ausnahme (\cref{sec:schnittstellen}). Tritt
während eines Zyklus eine Ausnahme auf, gleich in welchem Modul, wechselt der Orchestrator in
den Zustand \texttt{ERROR}, hält die Fehlermeldung zusammen mit dem Zustand fest, in dem sie
auftrat, und stoppt den Roboter. Ein Dauerbetrieb wird dabei beendet. Weiter geht es erst,
wenn der Bediener den Fehler quittiert: Dann wird der Roboter erneut vorbereitet
(\cref{sec:m2-api}), fährt in die Ruheposition, und der Automat kehrt nach \texttt{IDLE}
zurück.

Zusätzlich gibt es in der Oberfläche einen Stopp-Knopf, der den Roboter unabhängig vom
Zustand des Ablaufs sofort anhält. Er läuft in einem eigenen Thread, damit er nicht hinter
einer gerade laufenden Bewegung in der Warteschlange des Roboter-Threads wartet. Der
Software-Stopp ergänzt den Hardware-Not-Halt, ersetzt ihn aber nicht.

Bevor der Roboter an der realen Anlage aus der Oberfläche heraus bewegt wird, fragt ein Dialog,
ob der Arbeitsraum frei ist und der Not-Halt in Reichweite. Im Dauerbetrieb kann der Ablauf
außerdem so angehalten werden, dass der laufende Zyklus noch zu Ende geführt wird.

% ------------------------------------------------------------------------------
\section{Protokollierung}
\label{sec:int-protokoll}

Jeder Zyklus wird als Zeile in einer \ac{csv}-Datei protokolliert; pro Tag entsteht eine
Datei. Eine Zeile enthält Zeitpunkt und Nummer des Zyklus, ob er erfolgreich war, bei einem
Fehler dessen Meldung, die Dauer, die erkannte Pose, die Prüfergebnisse (Seriennummer,
Lochdurchmesser, Kerbe), das Gesamturteil und die Gründe für ein Schlechtteil. Optional werden
zu jedem Zyklus die Zwischenbilder von Objekterkennung und Sichtprüfung mit eingezeichneten
Ergebnissen gespeichert, sodass sich einzelne Zyklen nachträglich nachvollziehen lassen.

Diese Protokolle sind die Datengrundlage für die Evaluation in \cref{ch:evaluation}.
